Le stage effectué au sein de GIAS Industries a été une expérience d'une richesse exceptionnelle, marquant une étape déterminante dans notre formation de Licence en Informatique de Gestion. Ce projet de conception et de réalisation d'une plateforme de gestion des stagiaires a permis de répondre à une problématique réelle de transformation digitale au sein d'une grande structure industrielle.

Au cours de ces quelques mois, nous avons pu valider nos compétences techniques en développement Full-stack (React, Node.js, Sequelize) tout en découvrant les subtilités de la gestion des ressources humaines. L'application est aujourd'hui fonctionnelle et offre à la DRH un outil centralisé, sécurisé et automatisé, remplaçant avantageusement les méthodes de suivi manuelles.

\vspace{1cm}

Sur le plan professionnel, cette immersion chez GIAS Industries nous a appris l'importance de l'esprit d'équipe, de la communication inter-départementale et de l'adaptation constante aux besoins des utilisateurs finaux. Nous avons également pris conscience des enjeux de sécurité et de confidentialité liés aux données personnelles des collaborateurs et stagiaires.

Sur le plan personnel, ce travail a renforcé notre autonomie et notre capacité à résoudre des problèmes complexes sous la supervision de professionnels exigeants. Les conseils prodigués par nos encadrants, M. Abdelbasset Trad et M. Hazem Belazi, ont été d'une valeur inestimable pour mener à bien ce projet.

\vspace{1cm}

Bien que les objectifs initiaux aient été atteints, plusieurs perspectives d'évolution peuvent être envisagées pour pérenniser et enrichir la plateforme :
\begin{itemize}
    \item L'intégration d'un module de signature électronique pour les conventions de stage.
    \item L'ajout d'une application mobile pour faciliter le pointage des présences sur le terrain.
    \item L'interfaçage avec le système de paie existant du groupe HMS pour une gestion encore plus intégrée.
    \item L'utilisation de l'intelligence artificielle pour trier et classer automatiquement les CV reçus.
\end{itemize}

En conclusion, ce stage a constitué une véritable symbiose entre le développement de nos compétences techniques et notre épanouissement personnel, renforçant notre maturité face aux exigences du monde industriel. Entre rigueur méthodologique et enrichissement humain, cette expérience nous a permis de structurer notre vision de carrière. Forts de ces acquis et tournés vers l'avenir, nous envisageons avec enthousiasme de poursuivre cette dynamique d'innovation en explorant les nouvelles technologies qui façonneront les outils de gestion de demain.